% ============================================================
% SECTION 4.3-4.5: THE NEGATIVE, ITS RIVAL EXPLANATIONS, AND THE EXCEPTION.
% Restructured 2026-07-26 (late) from sections/04_below.tex, whose full previous
% form is preserved at sections/_mined/04_below_v1_full.tex.
%
% Changes in this pass, all structural:
%   - promoted from a standalone section to subsections of Experiments
%   - the fixed-k result and the speed result become paragraphs of 4.3 rather
%     than subsections of their own, since each is one measurement
%   - the recovery subsection is deleted: it restated the com-DBLP anatomy that
%     4.2 now owns, so 4.3 points at Table 5 instead of retelling it
%   - "four ways the loss could be an artifact" keeps its own subsection because
%     its four legs need paragraph headings of their own
% No numbers changed in this pass.
% STYLE: no em-dashes, no AI-characteristic phrasing (see STYLE.md).
% ============================================================

\subsection{Where sparsification does not help}\label{sec:experiments:negative}

\paragraph{The objective}

Table~\ref{tab:negative} reports the result. Across 102 matched-retention cells,
seven sparsifiers on seven networks under a modularity optimizer, not one cell
improves modularity measured on the original graph, whether the baseline is the
compute-matched control or the best of five plain restarts. The best cell in the
entire grid loses $0.0105$.

\begin{table}[t]
\centering
\caption{Every matched-retention cell, Leiden, seven networks. $\Delta Q$ is
measured on the original graph, Eq.~\eqref{eq:transfer}, against the best of five
plain restarts. Counts are cells with a positive value. The arms differ in how
much they lose and not in whether they lose.}
\label{tab:negative}
\small
\setlength{\tabcolsep}{4pt}
\begin{tabular}{@{}lrrrr@{}}
\toprule
Arm & Cells & $\Delta Q > 0$ & Best cell & Worst cell\\
\midrule
K-Neighbor            & 18 & 0 & $-0.0105$ & $-0.667$\\
DSpar                 & 18 & 0 & $-0.0116$ & $-0.640$\\
L-Spar                &  9 & 0 & $-0.0117$ & $-0.407$\\
Backbone, random fill & 14 & 0 & $-0.0141$ & $-0.214$\\
Local Similarity      & 15 & 0 & $-0.0172$ & $-0.311$\\
Local Degree          & 14 & 0 & $-0.0288$ & $-0.278$\\
Backbone, Jaccard fill& 14 & 0 & $-0.0290$ & $-0.235$\\
\midrule
Total                 & 102 & \textbf{0} & $-0.0105$ & $-0.667$\\
\bottomrule
\end{tabular}
\end{table}

The sweep that isolates one method across its full retention range agrees. Local
L-Spar on the same seven networks loses between $0.014$ and
$0.404$ against a compute-matched baseline in every configuration, with
the loss growing as retention falls. Its three most aggressive configurations
cannot reach the retention they were given and run at the floors of
Table~\ref{tab:floors} instead.

\paragraph{Under a fixed-$k$ partitioner}
The results above concern
detectors that choose their own granularity. The condition under which
sparsification does help, established in this paper on synthetic benchmarks,
involves a detector that takes the number of communities as an input, so the
question is whether such a detector escapes the negative on these graphs.

It does not. Metis was run on all seven networks at the same operating points and
under the same protocol, with ten partitioner seeds per cell and the worst
sparsified run compared against the best baseline run. On the five networks with
average degree between $5.5$ and $10.7$, no cell is positive: 128 cells, seven
sparsifiers, and both ways of choosing $k$ on the two of those networks that
carry reference communities. The best cell loses $0.017$. Fixing
the community count removes the granularity artifact by construction and leaves
the loss exactly where it was.

The two conventions for choosing $k$ agree throughout, whether $k$ is taken from
the community count a modularity optimizer returns on the original graph or, on
the three networks with reference communities, from the number of those
communities. On com-Amazon the two values differ by a factor of sixteen, 306
against 5,000, and both give the same answer.

Above that range the picture changes, and Section~\ref{sec:experiments:boundary} takes it up:
on one of the two networks above it, four of the seven sparsifiers produce a gain
that survives the worst case over ten seeds. On the other, which is denser, none
does.

\paragraph{Recovery}
Agreement with reference communities behaves differently from the objective, and
the difference is instructive rather than encouraging. Raw agreement often rises,
and Table~\ref{tab:granularity} has already shown why: the sparsified partition is
finer, the measure rewards that, and the floor beneath it rises at the same time.
After both controls something still remains on two of the three labelled
networks, and Section~\ref{sec:experiments:exception} takes those up.

\paragraph{Speed}
The pipeline is not faster. Of the 63 configurations, four preserve quality in the
sense of coming within $0.005$ of the best-of-$N$ baseline, and the fastest of
those four runs at $0.66\times$ the speed of clustering the original graph
directly. Twelve configurations exceed $1.5\times$, and every one of them loses
modularity, by $0.015$ to $0.162$. Timing only the detection step and giving the
sparsifier away for free, the median speedup is $1.21\times$, $1.63\times$ and
$1.15\times$ by algorithm, well short of the factor of two that halving the edge
count suggests, because all three algorithms return more communities on the
sparsified graph and do more work per pass.

\subsection{Four ways the loss could be an artifact}\label{sec:experiments:artifact}

A negative result is only as good as the explanations it excludes. Four are
available, and each was tested rather than argued.

\paragraph{Fragmentation}
Sparsification breaks graphs apart, and the optimizer may be penalized for being
handed a shattered graph. The backbone arms test this directly, because they
retain a spanning structure and cannot increase the number of connected
components. They behave as designed: under the backbone with random fill no
detected community is a singleton in any of its fourteen cells, and the fraction
of vertices in communities of fewer than ten members never exceeds $1.8\%$,
where the shattering arms reach $86\%$ on the same networks. The conclusion does
not move. The backbone's best cell is $-0.0141$, worse than three of the arms that
fragment freely. Removing half the edges costs modularity on the original graph
whether or not the graph stays in one piece.

\paragraph{The choice of optimizer}
Everything above uses a modularity optimizer, so the objective and the algorithm
may be too closely coupled. Two sparsifiers across three further algorithms, one
flow-based, one a local majority rule and a second modularity optimizer, give the
63 configurations of Table~\ref{tab:accounting}. Against the best of $N$ plain
restarts, 60 are negative, with mean changes of $-0.082$ under Infomap, $-0.046$
under Louvain and $-0.106$ under label propagation. The three positive cells are
label propagation on one network whose unsparsified baseline is pathological: its
own best of twenty restarts reaches modularity $0.079$ where a modularity
optimizer on the same graph reaches $0.416$. Relieving a known failure mode of one
algorithm is not a benefit of sparsification.

\paragraph{A benign loss}
The loss might be harmless if sparsification sheds peripheral vertices while
leaving community cores intact. It does not. The fragments subdivide communities
without mixing them, but they cut the most embedded vertices as readily as the
least, and a resolution-matched partition of the original graph preserves cores
better in sixteen of twenty comparisons. Restricting the damage to leaf and
near-leaf vertices does not rescue it either.

\paragraph{The retention regime}
The loss might be an artifact of pushing the methods too hard. It is not: it is
already present at the mildest retention any of them can reach. On the three
sparsest networks that is about a third of the edges, by Table~\ref{tab:floors},
and every arm loses there. The regime in which sparsification would pay for itself
is not merely unprofitable on these graphs, it is unreachable.

One observation from the same arms does not belong to this list, because it
supports no rival explanation. The backbone's two fills differ only in which edges
top up the spanning structure, random or highest Jaccard similarity, and random
fill scores better in ten of the fourteen paired cells. In this regime the
similarity signal subtracts.

\subsection{The exception}\label{sec:experiments:exception}

Two networks carry a result that clears the matched control of
Section~\ref{sec:protocol:granularity} and the chance floor.

On com-Amazon, L-Spar at realized retention $0.542$
reaches average best-match $F_1$ of $0.402$ against a baseline of $0.142$. The
matched control does not remove it: partitioning the original graph to the
same community count gives $0.319$, and deliberately over-matching by 37 per cent
more communities still gives only $0.372$. The chance floor does not remove it
either, and subtracting the floor widens the margin rather than narrowing it,
because the floor rises with the community count. A second instance
appears in the same experiments from an unrelated method, Local Degree, on
com-Amazon and on com-DBLP. The two baselines separate on com-DBLP: that instance
clears the matched control, at $0.163$ against $0.123$ at the same community
count, and does not clear the sweep, whose best point on the untouched graph
reaches $0.397$.

What removes both is a different comparison. Both non-modularity algorithms reach
higher agreement on com-Amazon without any sparsification at all, $0.465$ under
Infomap and $0.480$ under label propagation, against the $0.402$ that the
sparsified modularity pipeline attains, and the same holds on com-DBLP at
$0.326$ and $0.382$ against $0.163$. The effect is real, and it is the
modularity objective's default granularity being partially repaired by shattering
the graph, in a metric that rewards fine partitions. It is not a capability that
sparsification confers.

Two things about this cell are worth carrying forward. It coexists with a
modularity loss of $0.079$ in the same configuration, so the objective and the
recovery measure move in opposite directions on the same partition of the same
graph. And it is the reason the negative result of this section is stated as it
is: no benefit on the objective anywhere, and no recovery benefit that survives
comparison against a better algorithm run on the untouched graph. Section
\ref{sec:discussion} takes up the dissociation as an open phenomenon.
